% sec-01: the approval, the chronology, and the four things unchanged.
% Figure 1 (mermaid-type flowchart), Table 1.
\section{What Changed on August 26, 2026}

The FDA approved Rasonque (daraxonrasib), Revolution Medicines' RAS(ON)
multi-selective inhibitor, as the first-in-class targeted therapy for metastatic
pancreatic cancer \cite{fdarasonque2026}. That is one sentence, and this section
is about the difference between what it does change and what a reader in a hurry
might assume it changes.

\subsection{The dated chronology}

Four artifacts, each public, each dated, and each citable without contacting this
company \cite{daraxstory2026}.

In June 2025 a review of forty pancreatic ductal adenocarcinoma meta-analyses
totaling over 400,000 words ranked a daraxonrasib combination as the lead funding
candidate for this company \cite{daraxident}. In August 2025 a ten-arm
quantitative systems pharmacology simulation across roughly 250 ordinary
differential equations returned median overall survival of 12.8 months against
5.4 for chemotherapy \cite{qspmetpancre}. In May 2026, nine months later,
RASolute 302 reported 13.2 months against 6.6 in the RAS G12 population of
previously treated metastatic pancreatic cancer \cite{rasolute302}. In August
2026 the agent was approved.

The simulated ratio was 2.4-fold and the observed ratio was 2.0-fold. That
proximity is a chronology observation and a hypothesis-supporting one. It is
\textbf{not} a validation claim, and three differences are material: 1000
simulated against 241 enrolled; a combination against a single agent; and KRAS
G12C against a primarily G12D and G12V population. Those three clauses have
appeared in the same paragraph as those numbers in every document this company
has sent a federal office, and they appear here for the same reason.

\begin{appfloat}[!tb]
\begin{appfig}[mermaid-type / flowchart with a dated transition]
% Three columns on a 5.6cm pitch.  The before column converges on the transition
% node; the after column diverges from it.  The unchanged panel hangs below the
% node on a dashed edge, so it reads as a qualification of the arrow.
% Node names are plain letters and digits: a name carrying a dot would be read
% as a node-and-anchor pair, and a name carrying a minus as a subtraction.
\node[mmgray,text width=30mm] (bone)   at (0,0)      {Agent investigational, Phase 3 ongoing};
\node[mmgray,text width=30mm] (btwo)   at (0,-1.55)  {Funder underwrites agent, device, workflow};
\node[mmgray,text width=30mm] (bthree) at (0,-3.10)  {The 2025 simulation is a hypothesis};
\node[mmdec,text width=25mm]  (fda)    at (5.6,-1.55){FDA approval\\Rasonque\\August 26, 2026};
\node[mmgoal,text width=30mm] (aone)   at (11.2,0)     {Agent approved and labeled, metastatic};
\node[mmmid,text width=30mm]  (atwo)   at (11.2,-1.55) {Funder underwrites device and workflow};
\node[mmsoft,text width=30mm] (athree) at (11.2,-3.10) {The 2025 simulation is a dated public call};
\foreach \n in {bone,btwo,bthree}{\draw[mmedge]  (\n.east) -- (fda.west);}
\foreach \n in {aone,atwo,athree}{\draw[mmedgeb] (fda.east) -- (\n.west);}
% Unchanged panel: four boxes on a 2.95cm pitch, framed and titled.
\node[mmsoft,text width=25mm,font=\tiny\sffamily] (uone)   at (2.40,-5.35)  {Perioperative use still investigational};
\node[mmsoft,text width=25mm,font=\tiny\sffamily] (utwo)   at (5.35,-5.35)  {An IND is still required};
\node[mmsoft,text width=25mm,font=\tiny\sffamily] (uthree) at (8.30,-5.35)  {No supply agreement exists};
\node[mmsoft,text width=25mm,font=\tiny\sffamily] (ufour)  at (11.25,-5.35) {Robotic configuration not specified};
\node[mmlanetitle,anchor=west] (utitle) at (0.60,-4.55) {Unchanged by the approval};
\node[mmlane,fit={(uone)(ufour)(utitle)},inner sep=7pt] {};
\draw[mmedged] (fda.south) -- (5.35,-4.28);
\end{appfig}
\vspace{-0.60cm}
\figcaption{Figure 1. What the FDA approval changed in the ask, as three items moving\\
across one dated transition, and four items the approval leaves untouched.}
\end{appfloat}

\subsection{The four rows that matter}

\begin{apptable}
\begin{tabularx}{\textwidth}{>{\raggedright\arraybackslash}p{2.9cm} >{\raggedright\arraybackslash}p{4.6cm} >{\raggedright\arraybackslash}X}
\headrow \thead{Dimension} & \thead{Before the approval} & \thead{After the approval}\\
\midrule
Agent status & Investigational, Phase 3 evaluation ongoing & Approved and labeled in metastatic disease \\
Risk underwritten & Agent risk, device risk, workflow risk & Device risk and workflow risk \\
Supply question & Whether an investigational agent can be obtained at all & Whether a perioperative investigational use of an approved agent can be supported \\
Method evidence & A method argued from its own outputs & A method with one external checkpoint a reviewer can date \\
\end{tabularx}
\end{apptable}
\vspace{-0.60cm}
\tabcap{Table 1. The four dimensions the approval moves, stated as a before and\\
an after, with no fifth row claimed and no row overstated in either column.}

\subsection{The four things it does not change}

Each of the following is a place where an enthusiastic reading of the approval
would be wrong, and each is stated in every letter this packet accompanies.

The proposed use remains investigational. The approval covers metastatic disease;
this program proposes perioperative use in resectable and borderline-resectable
disease, and an investigational new drug application is still required
\cite{ecfr2026part312}.

No supply arrangement exists. No drug supply agreement, letter of authorization,
or regulatory cross-reference is in place with the agent's developer, and no
approach has been made \cite{revmedexternal2026}.

The device and software questions are untouched. The advisory boundary, the 3
millisecond arm-level stop, the 500 millisecond system-wide stop, the
verification harness, and the human-authority guarantee are exactly where they
were \cite{onpremwhippl}.

The novelty claim is unchanged. Daraxonrasib is nowhere described as first in
human. The supportable claim concerns the first prospective clinical evaluation
of the integrated surgical and advisory workflow, subject to FDA and
institutional confirmation.
